% GENERATED FILE. Edit canonical lessons, then run npm run generate:books.
% canonical-source-hash: fnv1a64:ce1af315168e1af1
% canonical-lessons: JA-C483-nendo, JA-C483-setomono, JA-C483-utsuwa, JA-C483-katameru, JA-C483-katamaru

\chapter{Clay and China}
\label{ch:ja-clay-and-china}
\hlchaptermodality{483}

\begin{chapteropening}
\textbf{By the end of this chapter:} \emph{I can talk about clay, pots and bowls.}

Complete the last lesson of chapter 483: I can talk about clay, pots and bowls.
\end{chapteropening}

\section[\texorpdfstring{\ja{ねんど} (nendo)}{ねんど (nendo)}]{\ja{ねんど} (nendo) — clay}

\label{lesson:JA-C483-nendo}



\begin{quote}
Before the new one: say the Japanese for made of wood, then the Japanese for to stick into.
\end{quote}

\subsection*{You'll want to know: \ja{ねんど}}
\textbf{\ja{ねんど}} — \emph{nendo} — \textquotedblleft{}clay\textquotedblright{}.

\begin{quote}
\textbf{\ja{こどもが} \ja{ねんどで} \ja{ねこを} \ja{つくりました。}} — \emph{kodomo ga nendo de neko o tsukurimashita} — \textquotedblleft{}The child made a cat out of clay.\textquotedblright{}
\end{quote}

\begin{cousinweb}
\textbf{\ja{ねんど}} (\emph{nendo}) is Sino-Japanese (kango): two Chinese characters read with their on'yomi, readings borrowed from Chinese. They are \emph{nen}, \textquotedblleft{}sticky\textquotedblright{} (Middle Chinese \emph{nrjem}) and \emph{do}, \textquotedblleft{}soil\textquotedblright{} (\emph{thuX}) (KANJIDIC2; Middle Chinese in Baxter's transcription).
\end{cousinweb}

\begin{grammarlens}[title={What you've built}]
One more word for describing what you lived and felt.
\end{grammarlens}

\subsection*{Your turn}
\begin{itemize}
\raggedright
  \item \emph{Say it:} \emph{nendo}
  \item \emph{Say it:} \emph{nendo}, once more
  \item \emph{Say it:} \emph{kodomo ga nendo de neko o tsukurimashita} — The child made a cat out of clay.
  \item \emph{Recall:} say the Japanese for made of wood, then the Japanese for to stick into, then say \emph{nendo} again
\end{itemize}

\begin{tcolorbox}[breakable,colback=teal!4,colframe=teal!35!black,title={Before you move on}]
What does \ja{ねんど} mean? (\textquotedblleft{}Clay\textquotedblright{}.) Say it once more.
\end{tcolorbox}
\section[\texorpdfstring{\ja{せともの} (setomono)}{せともの (setomono)}]{\ja{せともの} (setomono) — china; earthenware}

\label{lesson:JA-C483-setomono}



\begin{quote}
Before the new one: say the Japanese for to stick into, then the Japanese for clay.
\end{quote}

\subsection*{You'll want to know: \ja{せともの}}
\textbf{\ja{せともの}} — \emph{setomono} — \textquotedblleft{}china; earthenware\textquotedblright{}.

\begin{quote}
\textbf{\ja{せとものの} \ja{ちゃわんを} \ja{かいました。}} — \emph{setomono no chawan o kaimashita} — \textquotedblleft{}I bought a china teacup.\textquotedblright{}
\end{quote}

\begin{cousinweb}
\textbf{\ja{せともの}} (\emph{setomono}) is a native Japanese compound (wago), inherited rather than borrowed from Chinese: \emph{seto}, written with the characters for \textquotedblleft{}shallows\textquotedblright{} and \textquotedblleft{}door\textquotedblright{}, and \emph{mono}, \textquotedblleft{}thing\textquotedblright{} (JMdict, KANJIDIC2).
\end{cousinweb}

\begin{grammarlens}[title={What you've built}]
One more word for describing what you lived and felt.
\end{grammarlens}

\subsection*{Your turn}
\begin{itemize}
\raggedright
  \item \emph{Say it:} \emph{setomono}
  \item \emph{Say it:} \emph{setomono}, once more
  \item \emph{Say it:} \emph{setomono no chawan o kaimashita} — I bought a china teacup.
  \item \emph{Recall:} say the Japanese for to stick into, then the Japanese for clay, then say \emph{setomono} again
\end{itemize}

\begin{tcolorbox}[breakable,colback=teal!4,colframe=teal!35!black,title={Before you move on}]
What does \ja{せともの} mean? (\textquotedblleft{}China; earthenware\textquotedblright{}.) Say it once more.
\end{tcolorbox}
\section[\texorpdfstring{\ja{うつわ} (utsuwa)}{うつわ (utsuwa)}]{\ja{うつわ} (utsuwa) — a bowl; a vessel}

\label{lesson:JA-C483-utsuwa}



\begin{quote}
Before the new one: say the Japanese for clay, then the Japanese for china.
\end{quote}

\subsection*{You'll want to know: \ja{うつわ}}
\textbf{\ja{うつわ}} — \emph{utsuwa} — \textquotedblleft{}a bowl; a vessel\textquotedblright{}.

\begin{quote}
\textbf{\ja{りょうりを} \ja{しろい} \ja{うつわに} \ja{もります。}} — \emph{ryōri o shiroi utsuwa ni morimasu} — \textquotedblleft{}I serve the food in a white bowl.\textquotedblright{}
\end{quote}

\begin{cousinweb}
\textbf{\ja{うつわ}} (\emph{utsuwa}) is a native Japanese word (wago), inherited rather than borrowed from Chinese; its kanji was attached later as a kun'yomi, and its older origin is not settled (JMdict, KANJIDIC2).
\end{cousinweb}

\begin{grammarlens}[title={What you've built}]
One more word for describing what you lived and felt.
\end{grammarlens}

\subsection*{Your turn}
\begin{itemize}
\raggedright
  \item \emph{Say it:} \emph{utsuwa}
  \item \emph{Say it:} \emph{utsuwa}, once more
  \item \emph{Say it:} \emph{ryōri o shiroi utsuwa ni morimasu} — I serve the food in a white bowl.
  \item \emph{Recall:} say the Japanese for clay, then the Japanese for china, then say \emph{utsuwa} again
\end{itemize}

\begin{tcolorbox}[breakable,colback=teal!4,colframe=teal!35!black,title={Before you move on}]
What does \ja{うつわ} mean? (\textquotedblleft{}A bowl; a vessel\textquotedblright{}.) Say it once more.
\end{tcolorbox}
\section[\texorpdfstring{\ja{かためる} (katameru)}{かためる (katameru)}]{\ja{かためる} (katameru) — to harden; to make firm}

\label{lesson:JA-C483-katameru}



\begin{quote}
Before the new one: say the Japanese for china, then the Japanese for a bowl.
\end{quote}

\subsection*{You'll want to know: \ja{かためる}}
\textbf{\ja{かためる}} — \emph{katameru} — \textquotedblleft{}to harden; to make firm\textquotedblright{}.

\begin{quote}
\textbf{\ja{ねんどを} \ja{てで} \ja{かためます。}} — \emph{nendo o te de katamemasu} — \textquotedblleft{}I press the clay firm with my hands.\textquotedblright{}
\end{quote}

Polite \textbf{\ja{かためます}}, \emph{katamemasu}; an ichidan verb. Past (polite): \textbf{\ja{こどもが} \ja{ゆきを} \ja{かためました。}}, \emph{kodomo ga yuki o katamemashita}.

\begin{cousinweb}
\textbf{\ja{かためる}} (\emph{katameru}) is a native Japanese word (wago), inherited rather than borrowed from Chinese, the transitive partner of \emph{katamaru}, \textquotedblleft{}to harden\textquotedblright{}; its kanji was attached later as a kun'yomi (JMdict, KANJIDIC2).
\end{cousinweb}

\begin{grammarlens}[title={What you've built}]
One more word for describing what you lived and felt.
\end{grammarlens}

\subsection*{Your turn}
\begin{itemize}
\raggedright
  \item \emph{Say it:} \emph{katameru}
  \item \emph{Say it:} \emph{katameru}, once more
  \item \emph{Say it:} \emph{nendo o te de katamemasu} — I press the clay firm with my hands.
  \item \emph{Recall:} say the Japanese for china, then the Japanese for a bowl, then say \emph{katameru} again
\end{itemize}

\begin{tcolorbox}[breakable,colback=teal!4,colframe=teal!35!black,title={Before you move on}]
What does \ja{かためる} mean? (\textquotedblleft{}To harden; to make firm\textquotedblright{}.) Say it once more.
\end{tcolorbox}
\section[\texorpdfstring{\ja{かたまる} (katamaru)}{かたまる (katamaru)}]{\ja{かたまる} (katamaru) — to set hard; to harden}

\label{lesson:JA-C483-katamaru}



\begin{quote}
Before the new one: say the Japanese for a bowl, then the Japanese for to harden.
\end{quote}

\subsection*{You'll want to know: \ja{かたまる}}
\textbf{\ja{かたまる}} — \emph{katamaru} — \textquotedblleft{}to set hard; to harden\textquotedblright{}.

\begin{quote}
\textbf{\ja{ねんどは} \ja{いちにちで} \ja{かたまります。}} — \emph{nendo wa ichinichi de katamarimasu} — \textquotedblleft{}The clay sets in a day.\textquotedblright{}
\end{quote}

Polite \textbf{\ja{かたまります}}, \emph{katamarimasu}; a godan verb. Past (polite): \textbf{\ja{ねんどが} \ja{かたまりました。}}, \emph{nendo ga katamarimashita}.

\begin{cousinweb}
\textbf{\ja{かたまる}} (\emph{katamaru}) is a native Japanese word (wago), inherited rather than borrowed from Chinese, the intransitive partner of \emph{katameru}, \textquotedblleft{}to make firm\textquotedblright{}; its kanji was attached later as a kun'yomi (JMdict, KANJIDIC2).
\end{cousinweb}

\begin{grammarlens}[title={What you've built}]
One more word for describing what you lived and felt.
\end{grammarlens}

\subsection*{Your turn}
\begin{itemize}
\raggedright
  \item \emph{Say it:} \emph{katamaru}
  \item \emph{Say it:} \emph{katamaru}, once more
  \item \emph{Say it:} \emph{nendo wa ichinichi de katamarimasu} — The clay sets in a day.
  \item \emph{Recall:} say the Japanese for a bowl, then the Japanese for to harden, then say \emph{katamaru} again
\end{itemize}

\begin{tcolorbox}[breakable,colback=teal!4,colframe=teal!35!black,title={Before you move on}]
What does \ja{かたまる} mean? (\textquotedblleft{}To set hard; to harden\textquotedblright{}.) Say it once more.
\end{tcolorbox}
